%!TEX root = ../main.tex

\chapter{Bevezetés}


\section{Motiváció}

Az algoritmikus zeneszerzés régi terület: Hiller és Isaacson \cite{hiller1958musical}-ja, ahol egy IBM gépen sztochasztikus szabályrendszer írta meg a négy tételt. Évtizedeken át a kérdés úgy szólt: \emph{milyen szabályrendszer} írja le a tonális zenét. A mélytanulás megjelenésével a kérdés átalakult: ahelyett, hogy formális zeneelméleti szabályokat kódolnánk, \emph{milyen mintázatokat} talál egy hálózat egy meglévő zenei korpuszban; a terület átfogó áttekintéseit lásd \cite{DBLP:journals/corr/abs-1709-01620} és \cite{ji2023comprehensive} munkáiban. Ez a dolgozat a második szempontra koncentrál, méghozzá szimbolikus (MIDI) reprezentáció, ahol az események diszkrétek és pontosan megcímkézhetők.

A puszta mintázattanulás azonban önmagában nem elég, mert egyetlen modell sem tudja egyszerre lefedni a zenegenerálás összes aspektusát: a lokális dallami kontúrt, a transzpozíció-invariáns stílust és a hosszabb távú szerkezetet. Ez a tapasztalat motiválta a rendszer tervezését, amelyben három különböző modellt osztok szét három különböző zenei aspektus között, és nem egyetlen modellből próbálok mindent kihozni.

Ebben a kutatásban azt vizsgálom, hogyan lehet három különböző adaptív módszert egymás mellé tenni úgy, hogy mindegyik a zenei alkotás más oldalát kezelje. A Markov-láncot (rejtett Markov-modell kiterjesztéssel) a helyi, skálához kötött dallami statisztikára használom: ez a rész be tudja olvasni a \cite{CuthbertAriza2010}-en keresztül a hangnemet, és csak az ahhoz tartozó hangokból mintavételez. A VAE (és annak GOLC-kiterjesztése) egy 128-dimenziós pitch-eloszlás látens manipulációját végzi: stílusbeli ,,arculatváltást'' itt lehet a legkönnyebben paraméterezni. A Transformer pedig a többsávos, polifón kíséretet (basszus, dobok, akkordok) generálja cross-attention koordinációval. Nem azt állítom, hogy ez az osztás egyedüli helyes felbontás, csak hogy a három modellnek így van a legkevésbé átfedő felelőssége.

\section{Kutatási célok}

A dolgozat elsődleges célja a rendszer kidolgozása és elemzése: ez egy hibrid mélytanulási architektúra szimbolikus zenegenerálásra. Három egymással összefüggő kérdést vizsgálok:

Először is, hogyan érdemes felosztani a generálási feladatot három különböző modell között úgy, hogy mindegyik saját, jól körülhatárolt zenei aspektust kezeljen, és a felelősségek minél kevésbé legyenek átfedők. A Zha-ban a Markov-lánc a hangnemtudatos dallami alapot adja, a VAE (és annak GOLC-kiterjesztése) a transzpozíció-invariáns látens manipulációt, a Transformer pedig a polifón kíséret és többsávos struktúra szintézisét.

Másodszor, milyen gyakorlati implementációs kihívások merülnek fel a három modell egymás melletti tanításánál és inferenciájánál --- a számítási költség, a tanítási stabilitás és a memóriakezelés szempontjából --- és milyen tervezési döntésekkel kezelhetők. Harmadszor, milyen formában építhető be zeneelméleti tudás (skála, akkordfunkció, hangnem-felismerés) az adatvezérelt modellekbe anélkül, hogy túlzottan korlátozná a kreatív kimenetet.

A dolgozat ezeket a kérdéseket részletes modell-specifikus fejezetekben (Markov, VAE/GOLC-VAE, Transformer) és egy integrációs fejezetben tárgyalja, amely a pipeline egészét írja le. A munka egyaránt szolgál gyakorlati implementációs leírásként és kísérleti platformként, ahol a hibrid generatív megközelítés lehetőségeit és korlátait elemezni lehet.